\section{Implementierung}
\label{sec:implementierung}

\subsection{Projektstruktur}

\begin{hinweisbox}[Platzhalter]
Verzeichnisbaum mit kurzer Erklärung jeder Datei, z.\,B.:
\begin{itemize}
    \item \texttt{FNN.java} Klasse \texttt{FNN}
\end{itemize}
\end{hinweisbox}

\subsection{Die Klasse \texttt{FNN}}

\begin{hinweisbox}[Platzhalter]
Beschreibung der Klasse: Welche Attribute (Liste von Gewichtsmatrizen, Bias-Vektoren, Layer-Größen, Lernrate), welche Methoden?
\end{hinweisbox}

\subsection{Forward Pass}

\begin{lstlisting}[language=Java, caption={Forward Pass in Java (schematisch)}, label={lst:forward}]
// Hier Code Platzieren
\end{lstlisting}

\subsection{Backpropagation}

\begin{lstlisting}[language=Java, caption={Backpropagation (schematisch)}, label={lst:backprop}]
// Hier Code Platzieren
\end{lstlisting}

\subsection{Hyperparameter}

\begin{table}[htbp]
    \centering
    \small
    \begin{tabular}{l l}
    \toprule
    \textbf{Hyperparameter} & \textbf{Wert} \\
    \midrule
    Anzahl Hidden Layers  & [HIER EINFÜGEN] \\
    Neuronen pro Layer    & [HIER EINFÜGEN] \\
    Aktivierungsfunktion  & [HIER EINFÜGEN] \\
    Lernrate $\eta$       & [HIER EINFÜGEN] \\
    Batch-Größe $B$       & [HIER EINFÜGEN] \\
    Anzahl Epochen $E$    & [HIER EINFÜGEN] \\
    \bottomrule
    \end{tabular}
    \caption{Verwendete Hyperparameter.}
    \label{tab:hyperparameter}
\end{table}

\newpage